\clearpage
\section*{Assessment of the challenge statement and of the AI-generated solution}

\paragraph*{Scope, provenance, and method}
The files assessed are those supplied for Challenge 56: \texttt{problem.tex}, the AI-generated \texttt{solution.tex}, \texttt{final\_answer.tex}, \texttt{numerics.py}, \texttt{answer.py}, \texttt{metadata.json} and \texttt{agent\_log.md}. None of them was modified. The AI solution and its script are reproduced verbatim in Appendices~A and~B. The solution was originally generated by \texttt{codex-cli/gpt-5.6-sol} at reasoning effort ``high'' on 2026-08-15 and is marked ``generated-unreviewed''. The agent had web access through a server-side search tool (21 searches, queries logged, returned pages not) and a Python shell without network access (3 commands, one scratch script that the agent deleted before the run ended). Its citations are bare URLs. Its \texttt{numerics.py} runs (Python~3.11) and reproduces the twelve-digit values of its text, $3.01239295648\times10^{-23}$ and $1.44734540662\times10^{-24}\,\mathrm{GeV}^{-1}$, and hence its box.

The AI solution is assessed against the reference solution above. The AI-generated results are all reproduced exactly by adopting the AI's own model in an independent script, so that every discrepancy below is attributed to a named modeling choice and not to arithmetic.

\paragraph*{The challenge statement}
The statement is well posed in its core: a single scalar coupling, a given field density, a given interferometer geometry and noise, a given frequency and two observation times determine the threshold coupling by an unambiguous chain (Steps~1 to~7 above). Four inputs that the statement does not fix nevertheless enter the answer beyond the second significant figure.
\begin{enumerate}[label=(\roman*),leftmargin=*]
\item The beamsplitter material and the laser wavelength. Only the refractive index $n=3.5$ is given. It identifies silicon at $2\,\mu$m, the Cosmic Explorer Stage~2 design, but this is an inference. The material enters the index-dispersion response $\kappa_n$, which is $4.4\%$ of the signal coefficient (Step~3), and the relativistic correction $K_\alpha$, about $1\%$.
\item The meaning of the single velocity $230\,\mathrm{km\,s^{-1}}$. Read as the one-dimensional velocity dispersion it fixes the coherence time $\tau_c$. If it is instead the circular speed $v_0$ of the standard halo model, $\tau_c$ doubles and the $0.7$~yr answer moves by $10\%$ (Step~6 and Table~\ref{tab:results}).
\item The detection statistic beyond the coherence time. The statement says nothing. The sharp two-regime rule used throughout the literature and the finite-coherence statistic $R(T)$ of Step~6 differ by $3\%$ at $1000$~s and $23\%$ at $0.7$~yr.
\item The definition of ``SNR $=1$'' (amplitude or power, expectation or exclusion). Both solutions use the amplitude convention $h_0\sqrt{T/S_h}=1$. A confidence-level exclusion would add a factor of about $3$ \citep{Centers2021}.
\end{enumerate}
The statement therefore determines the answer to two significant figures at $1000$~s and to the first figure with a $30\%$ band at $0.7$~yr. It cannot be corrected unambiguously (each missing input has a most plausible value but not a unique one), so \texttt{problem.tex} is left untouched and the assumptions are stated in the solution. The overdensity factor $178$, the $6$~cm thickness and the $40$~km arm length are consistent with the Cosmic Explorer design and raise no issue.

\paragraph*{The AI solution, item by item}
Table~\ref{tab:ai56} lists every modeling element of the AI solution against the reference solution, with a classification along two axes. The first axis separates hidden curriculum (an expert input that the statement does not supply and that a solver cannot derive from it), a genuine error of the solver (the AI), and a deficiency of the challenge statement itself (the question prompt). The second gives the severity: consequential if the boxed answer moves beyond the convention band of the statement, minor if it does not, and simplification if a known effect was dropped, stated or unstated. Both axes are recorded in the ``Classification'' column. The ``Effect'' column gives the quantitative change of the AI answer from that item alone, on which the severity is judged.

\begingroup\footnotesize\setlength{\tabcolsep}{4pt}
\begin{longtable}{p{3.0cm}p{3.1cm}p{3.4cm}p{4.2cm}p{1.9cm}}
\caption{The AI solution for Challenge 56 against the reference solution. ``Effect'' is the change of the AI's threshold $\Lambda_\gamma^{-1}$ relative to the reference solution from that item alone (negative means the AI claims a better sensitivity).}\label{tab:ai56}\\
\toprule Item & AI solution & Reference & Classification & Effect \\ \midrule\midrule
\endfirsthead
\multicolumn{5}{l}{\textit{Table \ref{tab:ai56}, continued}}\\ \toprule Item & AI solution & Reference & Classification & Effect \\ \midrule\midrule
\endhead
\midrule \multicolumn{5}{r}{\textit{continued on next page}}\\
\endfoot
\bottomrule
\endlastfoot
Field amplitude $\phi_0=\sqrt{2\rho}/m_\phi$; $m_\phi c^2=hf$ & $39.9907$~GeV; $8.2713\times10^{-22}$~GeV & $39.991$~GeV; same & correct, same conventions & none \\
\midrule
Beamsplitter coefficient & $n\,l=0.2100$~m, from \citet{Miller2022}, Eq.~(3) & $n\cos\theta_r\,l$ \newline $+\,l\,\partial_n(n\cos\theta_r)\,\kappa_n n$ \newline $=0.2152$~m (Steps~3 and~4) & simplification, stated only through the citation. The cited formula is a secondary transcription of \citet{GroteStadnik2019}, Eq.~(17), whose thickness term is $\sqrt2\,n\,l$, and predates the correction $\sqrt2(n-\tfrac12)l$ of \citet{Vermeulen2021}, Eq.~(5). Neither primary formula was cross-checked, although the agent log shows the $\sqrt2$ form was searched for. The AI value is within $2\%$ of the exact thickness term $n\cos\theta_r\,l$ by cancellation, not by derivation: the primary formulas would have given $-28\%$ and $-16\%$. & $+2.4\%$ \\
\midrule
Index-dispersion term $\kappa_n$ & omitted, not mentioned & $-0.0442$ from the silicon dispersion at $2\,\mu$m & prompt deficiency (material and wavelength not given); hidden curriculum & $+4.4\%$ (in the row above) \\
\midrule
Snell refraction, sign structure of the plate & absent & $\theta_r=11.66^\circ$, per-encounter $n\cos\theta_r$ & hidden curriculum; minor & $-2.1\%$ (in the row above) \\
\midrule
Finite light-travel time of the $40$~km arms & absent & $\cos(\omega L/c)/\mathrm{sinc}(\omega L/c)=0.9906$ & hidden curriculum; minor & $-0.94\%$ \\
\midrule
Relativistic $K_\alpha$, radiative $m_e$ shift & absent & about $1\%$ each, convention-dependent & convention; minor & about $1\%$ \\
\midrule
Coherence time & $\tau_c=1/[f(v/c)^2]=8494.85$~s & $\tau_{\rm coh}=2\pi/(m_\phi v^2)=8494.9$~s (same); Derevianko $\tau_c=\tau_{\rm coh}/2\pi=1352$~s used inside $R(T)$ & consistent; the log shows the $2\pi$ ambiguity was recognized & none by itself \\
\midrule
Statistic beyond $\tau_c$ & sharp two-regime rule, $\sqrt{S_h/T}$ then $\sqrt{S_h}/(T\tau_c)^{1/4}$ & finite-coherence $R(T)$: $1.029$ at $1000$~s, $9.24$ at $0.7$~yr against $7.14$ & hidden curriculum and convention; consequential at $0.7$~yr & $-2.8\%$ ($1000$~s) \newline $-22.7\%$ ($0.7$~yr) \\
\midrule
SNR convention & amplitude, $h_0\sqrt{T/S_h}=1$ & same & consistent & none \\
\midrule
Constants, units, Julian year & correct ($1\,\mathrm{cm}=5.0677\times10^{13}\,\mathrm{GeV}^{-1}$, $h=4.135667696\times10^{-24}$~GeV~s) & same & correct & none \\
\midrule
Precision & twelve digits in the text, six in the box, for inputs given to one to three & three, with the convention band stated & minor, presentation & none \\
\midrule
Assumptions and idealizations & none stated beyond the photon-only coupling (no gravity, uniform field, arm cavities, velocity meaning) & Steps~1, 5, 6 and the two closing paragraphs & documentation & none \\
\midrule
Net & $3.012\times10^{-23}$, $1.447\times10^{-24}$ & $3.05\times10^{-23}$, $1.85\times10^{-24}$ & & $-1.3\%$, $-22\%$ \\
\end{longtable}
\endgroup

The AI solution is, in short, the sharp two-regime estimate of the literature with the simplest published signal coefficient. Against the two-regime variant of the reference solution ($2.97\times10^{-23}$ and $1.43\times10^{-24}$, Table~\ref{tab:results}) it is $1.5\%$ high at both observation times, which is the sum of the beamsplitter-coefficient and light-travel rows of Table~\ref{tab:ai56}. Against the reference solution itself it is $1.3\%$ low at $1000$~s and $22\%$ low at $0.7$~yr, the latter entirely from the statistic. No step is a hallucination: every formula it uses is published, and the one citation that could not be identified from its URL alone, \href{https://doi.org/10.1103/PhysRevD.105.103035}{10.1103/PhysRevD.105.103035}, is \citet{Miller2022}, whose Eq.~(3) does contain the coefficient $n\,l$.

\paragraph*{The AI numerical script}
\texttt{numerics.py} (Appendix~B) is under forty lines and mirrors the text line by line with exact CODATA constants, and it asserts agreement with \texttt{answer.py} to $5\times10^{-6}$. It runs and passes. It has no inputs block (every value is a literal inside one function), no sanity check of any intermediate quantity, no variant, and no comment on conventions. Adopting its model in an independent script reproduces its two numbers to ten digits, the residual being its truncated Planck constant. As numerical work it is correct and minimal.

\paragraph*{Evidence from the agent log}
The log records three points that the delivered text does not. First, the agent searched for the primary form of the coefficient (queries \texttt{"delta (L\_x-L\_y)" "sqrt\{2\}" beamsplitter} and \texttt{"n l" "Lambda\_gamma" beamsplitter}) and noted ``Evaluating geometric correction impact'' and ``Analyzing refractive index and approximation errors'', so the $\sqrt2$ and the index question were seen and then dropped without a word in the text. Second, it wrote ``I'm checking the precise $2\pi$ convention now, since it materially affects the $0.7$-year number'', so the coherence-time ambiguity was recognized, and the standard choice was made silently. Third, its scratch script printed a symbolic threshold \texttt{sqrt(2)*L*hthr*mass/(2*ell*n*sqrt(rho))}, which is the $n\,l$ model, and the script was then deleted, so the intermediate values are recoverable only from the log. The pattern is a solver that identified the right questions and documented none of them.

\paragraph*{Verdict}
The solution text is within expectation for a fast expert sketch: the chain $\phi_0\to h_0\to$ coherence regime $\to$ threshold is right, the numbers are right within the model, the units are clean. It is weak on justification: the central coefficient is taken from a secondary source without cross-check, and none of the four dropped effects (index dispersion, refraction, light travel, $K_\alpha$) is acknowledged, nor is any assumption stated. The script is correct and minimal. Holistically the AI answer is right to two significant figures at $1000$~s and fails at the third through secondary effects of $1$ to $4\%$ each that partly cancel, and it is $22\%$ low at $0.7$~yr through a convention that the statement does not fix and that the field itself uses. There is no consequential error. On the three-figure precision the challenge asks for, the statement's own inputs do not determine the third figure, and this should be reflected in the rubric.

\paragraph*{Recommended reference values and tolerance for grading}
Reference: $\Lambda_\gamma^{-1}=3.05\times10^{-23}\,\mathrm{GeV}^{-1}$ at $1000$~s and $1.85\times10^{-24}\,\mathrm{GeV}^{-1}$ at $0.7$~yr (finite-coherence statistic). Accepted alternative: $2.97\times10^{-23}$ and $1.43\times10^{-24}\,\mathrm{GeV}^{-1}$ (sharp two-regime rule, the literature convention). A tolerance of $\pm5\%$ around either pair at $1000$~s covers the geometry, index and $K_\alpha$ conventions, and $\pm12\%$ at $0.7$~yr covers in addition the reading of the velocity. Answers built on the primary-source coefficients without refraction, $\sqrt2\,n\,l$ or $\sqrt2(n-\tfrac12)l$, fall $28\%$ or $16\%$ below the reference solution and should be graded as partially correct. Values off by more than a factor of two indicate a wrong field amplitude, a wrong unit conversion, or a fully coherent $0.7$~yr integration (which would give $2.0\times10^{-25}$, a factor $9$ too optimistic).

\clearpage
\section*{Appendix A: the AI-generated solution and final answer, verbatim}
\begingroup
\small
\subsubsection*{Solution}
\paragraph*{Step-by-Step Derivation}

The scalar field amplitude and beamsplitter differential-displacement response are \href{\detokenize{https://link.aps.org/accepted/10.1103/PhysRevD.105.103035}}{given in the interferometer dark-matter literature}:

\[
\phi(t)=\phi_0\cos(2\pi f t), \qquad \phi_0=\frac{\sqrt{2\rho_{\rm DM}}}{m_\phi}
\]

in natural units $\hbar=c=1$, and, when only the photon coupling is nonzero,

\[
\delta(L_x-L_y) \simeq \Lambda_\gamma^{-1}\,n l\,\phi_0\cos(2\pi f t).
\]

Thus the peak strain amplitude is

\[
h_0=\frac{\delta(L_x-L_y)}{L} =\Lambda_\gamma^{-1}\frac{nl}{L}\phi_0.
\]

The finite field coherence time and the sensitivity scalings used below agree with the single-interferometer treatment in \href{\detokenize{https://arxiv.org/pdf/1906.06193}}{Grote and Stadnik}: sensitivity scales as $T^{-1/2}$ for $T\lesssim\tau_c$, and as $(T\tau_c)^{-1/4}$ for $T\gtrsim\tau_c$.

\subparagraph*{1. Dark-matter density and mass}

The overdense density is

\[
\rho_{\rm DM}=178(0.4) =71.2\ {\rm GeV\,cm^{-3}}.
\]

Using

\[
1\ {\rm cm}=5.06773071616\times10^{13}\ {\rm GeV^{-1}},
\]

we obtain

\[
\rho_{\rm DM} = \frac{71.2}{(5.06773071616\times10^{13})^3} = 5.47065597038\times10^{-40}\ {\rm GeV^4}.
\]

At $f=200\,{\rm Hz}$,

\[
m_\phi c^2=hf,
\]

so in energy units,

\[
m_\phi = (4.135667696\times10^{-24}\ {\rm GeV\,s})(200\ {\rm s^{-1}}) = 8.27133539200\times10^{-22}\ {\rm GeV}.
\]

Therefore,

\[
\phi_0 = \frac{\sqrt{2\rho_{\rm DM}}}{m_\phi} = \frac{\sqrt{2(5.47065597038\times10^{-40})}} {8.27133539200\times10^{-22}} = 39.9907057704\ {\rm GeV}.
\]

\subparagraph*{2. Strain per unit coupling}

For

\[
n=3.5,\qquad l=0.06\ {\rm m},\qquad L=4.00\times10^4\ {\rm m},
\]

the strain amplitude is

\[
h_0 = \Lambda_\gamma^{-1} \left(\frac{3.5(0.06)}{4.00\times10^4}\right) (39.9907057704\ {\rm GeV}),
\]

or

\[
h_0 = \left(2.09951205295\times10^{-4}\ {\rm GeV}\right) \Lambda_\gamma^{-1}.
\]

Here $\Lambda_\gamma^{-1}$ is expressed in ${\rm GeV^{-1}}$.

\subparagraph*{3. Coherence time}

The virialized scalar field has a fractional linewidth of order $(v/c)^2$, giving

\[
\tau_c \simeq \frac{1}{f(v/c)^2}.
\]

With $v=230\,{\rm km\,s^{-1}}$,

\[
\frac{v}{c} = \frac{230}{299792.458} = 7.67197\times10^{-4},
\]

and hence

\[
\tau_c = \frac{1} {200(230/299792.458)^2} = 8494.85046065\ {\rm s}.
\]

For a one-sided strain amplitude spectral density

\[
\sqrt{S_h}=2.00\times10^{-25}\ {\rm Hz^{-1/2}},
\]

the unit-SNR strain threshold is

\[
h_{\min}=
\begin{cases}
\dfrac{\sqrt{S_h}}{\sqrt{T_{\rm obs}}},
& T_{\rm obs}\leq\tau_c,\\[8pt]
\dfrac{\sqrt{S_h}}{(T_{\rm obs}\tau_c)^{1/4}},
& T_{\rm obs}>\tau_c.
\end{cases}
\]

\subparagraph*{4. Observation for $1000\,{\rm s}$}

Since

\[
1000\ {\rm s}<\tau_c,
\]

the observation is coherent:

\[
h_{\min} = \frac{2.00\times10^{-25}}{\sqrt{1000}} = 6.32455532034\times10^{-27}.
\]

Equating $h_0=h_{\min}$ gives

\[
(\Lambda_\gamma^{-1})_{\min} = \frac{6.32455532034\times10^{-27}} {2.09951205295\times10^{-4}\ {\rm GeV}} = 3.01239295648\times10^{-23}\ {\rm GeV^{-1}}.
\]

\subparagraph*{5. Observation for $0.7\,{\rm yr}$}

Using $1\,{\rm yr}=365.25$ days,

\[
T_{\rm obs} = 0.7(365.25)(86400) = 2.2090320\times10^7\ {\rm s}.
\]

Because $T_{\rm obs}\gg\tau_c$, the field phases cannot be added coherently over the full observation:

\[
h_{\min} = \frac{2.00\times10^{-25}} {\left[(2.2090320\times10^7)(8494.85046065)\right]^{1/4}} = 3.03871912597\times10^{-28}.
\]

Consequently,

\[
(\Lambda_\gamma^{-1})_{\min} = \frac{3.03871912597\times10^{-28}} {2.09951205295\times10^{-4}\ {\rm GeV}} = 1.44734540662\times10^{-24}\ {\rm GeV^{-1}}.
\]

\subsubsection*{Final Answer}

\[
\boxed{
\left.\Lambda_\gamma^{-1}\right|_{(1000\,\mathrm{s},\,0.7\,\mathrm{yr})}
=\left(3.01239\times10^{-23},\,1.44735\times10^{-24}\right)
}
\]

\endgroup

\clearpage
\section*{Appendix B: the AI-generated numerical script, verbatim}
\lstinputlisting[style=answerstyle,language=Python,caption={\texttt{numerics.py} (AI-generated)}]{ai_numerics.py}
\lstinputlisting[style=answerstyle,language=Python,caption={\texttt{answer.py} (AI-generated)}]{ai_answer.py}

\clearpage
\section*{Appendix C: the numerical script of the reference solution}
\lstinputlisting[style=answerstyle,language=Python,basicstyle=\ttfamily\footnotesize,caption={\texttt{challenge56.py}}]{challenge56.py}
